\documentclass[11pt,a4paper]{article}

\usepackage[margin=1in]{geometry}
\usepackage[T1]{fontenc}
\usepackage[utf8]{inputenc}
\usepackage{lmodern}
\usepackage{hyperref}
\usepackage{xcolor}
\hypersetup{colorlinks=true,linkcolor=blue,urlcolor=blue,citecolor=blue}

\title{Approach and Evaluation of the LangGraph Agentic Application}
\author{}
\date{}

\begin{document}

\maketitle

\section{Approach}

The LangGraph agentic application was populated through an iterative synthesis process that began from the explicit prompt to design an agentic app grounded in the previously mined process artifacts rather than as a generic issue-handling workflow. In practical terms, this meant that the application was not specified top-down from software engineering intuition alone, but bottom-up from the object-centric and role-centric mining results produced in the earlier notebooks. The OCDFG models provided the main structural intuition that the \texttt{issue} object should be treated as the central coordination object around which state is maintained, decisions are made, and related artifacts are attached. The BPMN models then informed the stage-oriented shape of the graph, namely an intake phase, contextualization, intent classification, routing, guarded action selection, execution, and closure or waiting. The DECLARE-style constraints were reflected in the policy layer, especially in the use of approval gates, risk flags, and guarded execution rather than unconditional autonomous action. Finally, the textual role process descriptions were used to populate the role-specific agents themselves, because they gave the clearest account of how each role interprets incoming work, what kinds of decisions it makes, and what outputs it is expected to produce.

A central design choice in populating the app was that not every mined role was turned into a directly routable executable node. Instead, the final implementation selected the roles whose behavior was both sufficiently distinct and operationally useful in a GitHub issue workflow. This produced the role-specific agents \texttt{issue\_reporter\_agent}, \texttt{bot\_workflow\_agent}, \texttt{implementation\_agent}, \texttt{quality\_agent}, and \texttt{technical\_writer\_agent}. The \texttt{issue\_reporter\_agent} was specified as the clarification-oriented agent. It corresponds to the mined issue reporter behavior, which is centered on expressing a need, refining intent, supplying missing details, and making the issue actionable when the current report is incomplete or ambiguous. In the app, this agent is therefore invoked when the incoming issue appears underspecified, and its output is an analysis that requests clarification, reframes the problem, or identifies information gaps that must be resolved before implementation or review can proceed.

The \texttt{bot\_workflow\_agent} was specified from the mined bot-oriented process behavior, which emphasized routine, automatable, low-risk, and highly structured actions. In the app, this agent handles issues that are better interpreted as repository housekeeping, lightweight workflow triggers, or procedural tasks that can be reduced to repeatable actions. Its role is not deep semantic interpretation of product requirements, but conversion of recognized patterns into operational steps such as labeling, routing, simple follow-up, or other mechanizable repository actions. In this sense, it acts as the bridge between the mined automated behavior and the concrete GitHub tool layer exposed in the application.

The \texttt{implementation\_agent} is the abstraction that captures behavior mined from roles such as feature developer and, to some extent, implementation-facing contributor activity. Rather than exposing several narrowly separated development roles, the app consolidates them into a single agent responsible for interpreting feature requests, enhancements, and bug-fix type intents as implementation work. This agent was specified to analyze the issue in engineering terms, infer the likely code-facing task, and derive an actionable plan such as implementing a new capability, adjusting existing behavior, or preparing a development-oriented next step. The consolidation was deliberate: the mined roles showed similar implementation semantics, and the application required a practical executable role rather than a fully faithful organizational replica.

The \texttt{quality\_agent} was populated from the mined quality engineer behavior, where the recurring emphasis was on validation, correctness, reliability, and the review of whether a proposed or completed action satisfies expectations. In the app, this agent represents testing and verification reasoning. It is therefore the natural destination for issues framed as quality concerns, test gaps, regressions, or requests for validation-oriented actions. Its specification focuses less on producing new features and more on identifying checks, test scenarios, acceptance criteria, and evaluation logic that should accompany or precede execution.

The \texttt{technical\_writer\_agent} was specified from the mined technical writer role description, which centered on explanation, user-facing clarity, and the maintenance of documentation as a first-class process outcome. In the application, this agent handles issues that are best understood as documentation requests or communication-improvement tasks. Its output is an analysis of what documentation artifact should be updated, how the issue should be translated into explanatory content, and what concrete documentation action is warranted. This made it possible for the app to distinguish documentation work from implementation work even when both are described in natural language through GitHub issues.

These agents were then embedded into a LangGraph workflow with shared typed state, GitHub integration tools, prompt templates, and routing logic. The population process therefore combined mined process semantics with software architecture constraints: the mined artifacts determined what kinds of role behavior were meaningful, while the LangGraph implementation turned those behaviors into executable agent nodes with explicit interfaces, routable responsibilities, and bounded actions.

\section{Evaluation}

The evaluation of the application was conducted as a functional test on two real GitHub issues, referred to as issue \#1 and issue \#2. The purpose of this evaluation was not to benchmark predictive accuracy statistically, but to verify whether the app could interpret realistic issue text, route the issue to the role-specific agent that best matched the inferred intent, and derive a plausible corresponding action through LLM-based analysis. In this sense, the two issues served as scenario-based tests of end-to-end behavior: issue ingestion, context loading, intent interpretation, role routing, and action derivation.

Issue \#1 was a documentation-oriented issue. Its content was interpreted by the application as a request whose main objective was not changing system behavior or implementing code, but improving explanation or written project material. Based on this interpretation, the classifier and routing logic directed the issue to the \texttt{technical\_writer\_agent}. This routing was consistent with the role population described above: the issue exhibited the semantics of documentation maintenance rather than implementation or testing. The LLM-based analysis performed by the agent translated the user request into a documentation task, identifying that the relevant next action was to prepare or update explanatory material rather than produce code changes. The significance of this first test is that it demonstrated role separation in a meaningful way. The app did not simply reduce every issue to a generic software-development task; instead, it recognized that some issues belong to a documentation process and should therefore be handled by a role specialized in communication and documentation artifacts.

Issue \#2 functioned as a complementary test because it was intended to represent implementation-oriented work. In its earlier phrasing, the issue was initially routed incorrectly to the bot-related path, which exposed an important property of the system: the routing behavior depends strongly on how the issue text signals its intent. After the wording was made more explicit and included a clearer feature-oriented formulation, the application reinterpreted the issue as implementation work and routed it to the \texttt{implementation\_agent}. This outcome is important for two reasons. First, it confirmed that the underlying routing logic was sensitive to issue semantics rather than arbitrary issue identifiers. Second, it showed that once the issue text contained clearer evidence of enhancement or feature intent, the app could correctly derive the relevant work type. The LLM-based analysis of the implementation agent then produced a corresponding action in the form of a development-oriented recommendation or plan, namely that the issue should be treated as a feature implementation task rather than as documentation, clarification, or lightweight automation.

Taken together, the two issues provide a compact but useful functional evaluation of the app. Issue \#1 validated that documentation requests are recognized and routed to the technical writer specialization, where the resulting action remains documentation-centered. Issue \#2 validated that feature-oriented requests can be interpreted as implementation tasks and routed to the implementation specialization, where the resulting action becomes development-centered. At the same time, the behavior observed during the earlier phrasing of issue \#2 exposed an expected limitation of natural-language routing: the app is only as clear as the issue text it receives. This is not merely an implementation weakness, but also a realistic property of issue-driven workflows. Overall, the evaluation shows that the application can operationalize mined role knowledge in a live GitHub setting, using LLM-based interpretation to transform issue descriptions into role-appropriate analyses and next actions.

\end{document}

% Made with Bob
